% !PRINTED 570
\weberpage{570}
%<<CONT
and since $\sin\frac{11\pi}{34}$ is greater than $\sin\frac{3\pi}{34}$, the sum of these
two expressions, and with it the whole sum $(-1,\eta)$, is positive % [sic] the printed authority ends the sentence without punctuation

There results further, for the real part of $(i,\eta)$,

\[
\eta - \eta_2 + \eta_4 - \eta_6 = 2\cos\frac{2\pi}{17} + 2\cos\frac{\pi}{17} + 2\cos\frac{8\pi}{17} - 2\cos\frac{4\pi}{17},
\]

which is immediately recognised as positive; finally, for the real
part of $(\alpha,\eta)$,

\[
\eta - \eta_4 + \frac{1}{\sqrt{2}}(\eta_1 - \eta_3 - \eta_5 + \eta_7)
= 2\cos\frac{2\pi}{17} - 2\cos\frac{8\pi}{17} + \sqrt{2}\left(\cos\frac{6\pi}{17} + \cos\frac{3\pi}{17} + \cos\frac{7\pi}{17} - \cos\frac{5\pi}{17}\right),
\]

which likewise proves to be positive.

\begin{center}
\S.\,170.

\gesperrt{Properties of the Numbers} $\psi$.
\end{center}

The formulae for the resolvents derived in the preceding
section admit a wealth of applications, even without an exact
determination of the prime $n$, for special assumptions
about the number $e$; they thus give number-theoretic theorems for
whole categories of primes.

For algebra they provide arbitrarily many solutions of cyclic
equations of prescribed degree.

We precede the derivation of the most important of these applications
for algebra with some general considerations on
the functions $\psi$ which we defined in the preceding
section.

For the computation of the numbers $\psi_{\lambda,\mu}(\varepsilon)$ [\S.\,169, (5)] one makes
use, as we have seen in the case $n = 17$, of an index table.
But between these functions there exist relations by
which they can be reduced to a smaller number, on which
\gesperrt{Jacobi} has made remarkable investigations\footnote{Cf. \gesperrt{Kronecker}, Zur Theorie der \gesperrt{Abel'}schen Gleichungen. Journ. f. Mathem. Bd. 93.}.

The function $\psi_{\lambda,\mu}(\varepsilon)$ does not change, according to its definition,
when $\lambda,\mu$ are interchanged, although its expression
%CONT>>
